%% ma: L11-B14-0007
%% lop: 11
%% chuong: 4
%% bai: 14
%% dang: 11.14.3
%% loai: TN4
%% muc: TH
%% dap_an: "A"
%% nguon: "tai-lieu-hoc-tap-mon-toan-lop-11-theo-chuong-trinh-giao-duc-pho-thong-2018.pdf (Lớp Toán Thầy Duy, Phạm Lê Duy, THPT Chu Văn An), Chương 4 Bài 5, Phần 1, Câu 2 (tr.523)"
%% kien_thuc: "Lớp 11 Bài 14 (phép chiếu song song)"
%% trang_thai: cho-duyet
%% ly_do: "Dạng mới đề xuất (11.14.x), chờ giáo viên duyệt dạng; tài liệu gốc không có lời giải mẫu, đáp án do kho tự giải và kiểm bằng sympy."
%% ghi_chu: "Đáp án gốc A; lý luận của kho dựa vào tính chất phần trong của khối."
\begin{ex}
Khẳng định nào sau đây đúng?
\choice{\True Hình chiếu song song của một hình chóp cụt có thể là một hình tam giác.}{Hình chiếu song song của một hình chóp cụt có thể là một đoạn thẳng.}{Hình chiếu song song của một hình chóp cụt có thể là một hình chóp cụt.}{Hình chiếu song song của một hình chóp cụt có thể là một điểm.}
\loigiai{\sloppy Hình chóp cụt là khối có phần trong khác rỗng nên hình chiếu của nó lên mặt phẳng là một miền phẳng có phần trong khác rỗng: không thể là đoạn thẳng hay một điểm (B, D sai), và là hình phẳng nên không phải hình chóp cụt (C sai). Với chóp cụt tam giác, chiếu theo phương của đường thẳng nối đỉnh chóp ban đầu với một điểm của đáy lên mặt phẳng đáy, hai đáy chiếu thành hai tam giác lồng nhau, hình chiếu của cả khối là một tam giác (A đúng).}
\end{ex}
